\scanpage{18}
\uline{Interprétations en termes de filtrations du foncteur oubli \deleted{sur la catégorie}} des représentations linéaires de dimension finie de $G$. (Il s'agit de filtrations compatibles au produit tensoriel\inserted{et ``exactes',cette \unsure{dernière} cond. aut. si $\operatorname{car}.k=0$}). Supposons fixé une classe\inserted{$c$} de\inserted{conjugaison, rationelle sur $k$, de} telles filtrations (NB elle n'admet pas nécessairement un représentant rat/$k$); i\deleted{\unsure{o}}l\deleted{\illegible{1}} revient au même de se donner,\deleted{un} pour le Chevalley $G_o$ associé à $G$, un \deleted{homomorphisme} élément de $\operatorname{Hom}(\underline{G}_m,T_o)/W=M/W$ ($T_o=M\otimes\underline{G}_m$ est le tore maximal type de $G_o$) invariant par le sous-groupe\inserted{\handwritten{$\Gamma$}} de $\operatorname{Auext}G_o=E_o$ image de $\operatorname{Gal}(\bar{k}/k)$, ou encore un \deleted{homomorphisme $\underline{G}_m\unsure{\to}T$} élément de $M$ qui soit dans la chambre de Weyl de $M\otimes_{\underline{Z}}\underline{R}$ et qui soit stable par $\handwritten{\Gamma}$. On a alors le \unsure{schéma} $\handwritten{\underline{\varnothing}}_c$ des filtrations de classe $\handwritten{c}$, qui est un espace homogène sous $G$, de stabilisateur des points des paraboliques;\inserted{\deleted{\unsure{c'}}\handwritten{il} est\inserted{\handwritten{$\sim$}} un sous-schéma ouvert de $\uline{Par}(G)$, orbite rationelle sous $G$,\handwritten{soit $c'$,}} On peut regarder aussi le schéma $\operatorname{ST}_c$ des homomorphismes $i:\underline{G}_m\handwritten{\to}G$ de classe $c$, qui est donc une orbite rationelle de $G$ opérant sur $\uline{Homgr}(\underline{G}_m,G)$. I\deleted{\unsure{l|t}}l est canoniquement isomorphe au schéma des couples de Lévi $(P,L)$ tels que $P$ soit de classe $c'$. (En d'autres termes, grâce au choix d'un $c$, pour $P$\inserted{ou encore $\handwritten{\underline{\varnothing}}$} de classe $c'$ donné, il revient au même de se donner un $i:\underline{G}_m\handwritten{\to}P$ de la classe définie par $c$, où un sous-groupe de Lévi de $P$, à $i$ étant associé $\operatorname{Cent}(i)$.\handwritten{)} On a \deleted{un homo}\inserted{deux}morphismesde schémas canoniquement isomorphes
\[
\operatorname{ST}_c\handwritten{\longrightarrow\underline{\varnothing}}_c\quad,\quad
(\operatorname{CL})_{c'}\handwritten{\longrightarrow}\operatorname{Par}_{c'}\quad,
\]
le premier en termes de théorie des filtrations d\unsure{\deleted{u}a|\deleted{a}u} foncteurs fibres, le deuxième \deleted{d}en termes de groupes paraboliques.

Supposons maintenant que $k$ soit de car. nulle (j'ignore dans quelle mesure c'est vraiment nécessaire). On peut alors considérer le Schéma $\operatorname{Dyn}_c$ des homomorphismes $f:\operatorname{Sl}(2,k)\handwritten{\longrightarrow}G$ tels que $fi_o$ soit de classe $c$. Supposons-le non vide\inserted{\handwritten{et non réduit à un point,}} i.e. $c$ une classe de Dynkin\inserted{\handwritten{$\ne1$}}. Alors $\operatorname{Dyn}_c$ est un \deleted{\uline{torseur}}\inserted{\handwritten{espace homogène}} sous $G$\inserted{\handwritten{de stabilisateur $Z_o$ (il faudrait regarder la structure de ces $Z_o$; dimension, semi-simplicité).}}\deleted{(en fait, il est can. is. au torseur $\uline{Isom}(G_1,G)$, où $G_1$ est le groupe quasi-déployé correspondant à $G$).} Soit $\operatorname{Un}_c$ le schéma des éléments unipotents\inserted{\handwritten{de $G$}} dont la classe corresponde à \deleted{celle de} $c$ par la correspondance
% archive sheet number: 17
